\documentclass{article}
\usepackage{nips15submit_e,times}
\usepackage{hyperref}
\usepackage{url}
\usepackage{graphicx}
\usepackage{xcolor}
\usepackage{booktabs}
\usepackage{subcaption}
\usepackage{listings}
\usepackage{enumitem}

% List formatting
\setlist[itemize]{leftmargin=*, label=\textbullet, itemsep=0.2em, topsep=0.3em}
\setlist[enumerate]{leftmargin=*, itemsep=0.2em, topsep=0.3em}

% Code listing setup
\definecolor{codebg}{rgb}{0.95,0.95,0.95}
\definecolor{codekw}{rgb}{0.0,0.0,0.5}
\definecolor{codestr}{rgb}{0.5,0.0,0.0}
\definecolor{codecom}{rgb}{0.0,0.5,0.0}

\lstset{
    backgroundcolor=\color{codebg},
    basicstyle=\ttfamily\small,
    keywordstyle=\color{codekw}\bfseries,
    stringstyle=\color{codestr},
    commentstyle=\color{codecom}\itshape,
    frame=single,
    framerule=0.5pt,
    rulecolor=\color{gray},
    breaklines=true,
    breakatwhitespace=true,
    showstringspaces=false,
    numbers=left,
    numberstyle=\tiny\color{gray},
    captionpos=b,
    tabsize=4,
    language=Python
}

\title{\textbf{Psi -1 - TaskName Technical Writeup}}
\author{
Muhammad Aufar Rizqi Kusuma \\
Sekolah Teknik Elektro dan Informatika\\
Institut Teknologi Bandung\\
\texttt{13524061@std.stei.itb.ac.id} \\
\And
Kurt Mikhael Purba \\
Sekolah Teknik Elektro dan Informatika\\
Institut Teknologi Bandung\\
\texttt{13524065@std.stei.itb.ac.id} \\
\And
Mochamad Fachri Alfaridzi \\
Sekolah Teknik Elektro dan Informatika\\
Institut Teknologi Bandung\\
\texttt{13525128@std.stei.itb.ac.id} \\
}

\newcommand{\fix}{\marginpar{FIX}}
\newcommand{\new}{\marginpar{NEW}}

\nipsfinalcopy

\begin{document}

\maketitle

% ==================== PROBLEM OVERVIEW ====================
\section{Problem Overview}

 [Placeholder: Describe the task objective, evaluation metric, and what the model needs to predict/classify.]

The competition task requires us to [brief description of the problem]. The primary evaluation metric is [metric name], which measures [explanation of metric]. Our goal is to [what success looks like].

Key challenges identified:
\begin{itemize}
    \item Challenge 1: e.g., class imbalance, missing values, high dimensionality]
    \item Challenge 2: e.g., noisy labels, small dataset, computational constraints]
    \item Challenge 3: e.g., feature engineering complexity, time series dependencies]
\end{itemize}

% ==================== DATASET ====================
\section{Dataset}

 [Placeholder: Describe dataset structure, size, features, and any preprocessing applied.]

The dataset consists of [number] samples with [number] features. Key characteristics:

\begin{itemize}
    \item \textbf{Train set}: [number] rows, [number] columns
    \item \textbf{Test set}: [number] rows, [number] columns
    \item \textbf{Target variable}: [description of target, type, distribution]
    \item \textbf{Feature types}: [numerical, categorical, text, images, etc.]
\end{itemize}

Notable observations:
\begin{itemize}
    \item Observation 1: e.g., 15\% missing values in feature X]
    \item Observation 2: e.g., target distribution is skewed 80/20]
    \item Observation 3: e.g., high correlation between features A and B]
\end{itemize}

% ==================== EXPERIMENT ENVIRONMENT ====================
\section{Experiment Environment}

 [Placeholder: List hardware, software, libraries, and versions used.]

\begin{tabular}{ll}
    \toprule
    \textbf{Component} & \textbf{Specification}                          \\
    \midrule
    Platform           & Kaggle Notebook (GPU/TPU)                       \\
    Python Version     & 3.x                                             \\
    Key Libraries      & [e.g., scikit-learn 1.x, pandas 2.x, numpy 1.x] \\
    Deep Learning      & [e.g., PyTorch 2.x / TensorFlow 2.x / None]     \\
    Hardware           & [e.g., NVIDIA Tesla P100 / T4 / CPU only]       \\
    Execution Time     & [e.g., ~2 hours total]                          \\
    \bottomrule
\end{tabular}

% ==================== EDA ====================
\section{Exploratory Data Analysis (EDA)}

 [Placeholder: Summarize key findings from data exploration.]

Our EDA revealed several important patterns:

\begin{enumerate}
    \item \textbf{Target Distribution}: [describe distribution, any imbalance]
    \item \textbf{Feature Correlations}: [mention important correlations or lack thereof]
    \item \textbf{Missing Data Pattern}: [describe missingness, whether MCAR/MAR/MNAR]
    \item \textbf{Outliers}: [describe outlier detection and handling]
    \item \textbf{Interesting Patterns}: [any surprising findings that informed modeling]
\end{enumerate}

% Example figure placeholder:
% \begin{figure}[h]
%     \centering
%     \includegraphics[width=0.8\linewidth]{public/eda/target_distribution.png}
%     \caption{Target variable distribution}
%     \label{fig:target_dist}
% \end{figure}

% ==================== PIPELINE EXPLANATION ====================
\section{Pipeline Explanation}

 [Placeholder: Describe the full ML pipeline from preprocessing to prediction.]

\subsection{Preprocessing}

[Describe data cleaning, feature engineering, encoding, scaling, etc.]

\begin{lstlisting}[caption=Preprocessing snippet]
# Placeholder: Add your preprocessing code here
import pandas as pd
from sklearn.model_selection import train_test_split

df = pd.read_csv('train.csv')
# Your preprocessing steps
\end{lstlisting}

\subsection{Feature Engineering}

[Describe any created features, transformations, or selections.]

Key features created:
\begin{itemize}
    \item Feature 1: description and rationale
    \item Feature 2: description and rationale
    \item Feature 3: description and rationale
\end{itemize}

\subsection{Model Architecture}

[Describe the model(s) used, hyperparameters, and training strategy.]

\begin{itemize}
    \item \textbf{Model}: [e.g., LightGBM, XGBoost, Random Forest, Neural Network]
    \item \textbf{Hyperparameters}: [key parameters tuned]
    \item \textbf{Training Strategy}: [e.g., 5-fold CV, early stopping, learning rate schedule]
    \item \textbf{Ensemble}: [if applicable, describe blending/stacking approach]
\end{itemize}

\begin{lstlisting}[caption=Model training snippet]
# Placeholder: Add your model training code here
from sklearn.ensemble import GradientBoostingClassifier

model = GradientBoostingClassifier(
    n_estimators=100,
    max_depth=5,
    learning_rate=0.1
)
model.fit(X_train, y_train)
\end{lstlisting}

\subsection{Validation Strategy}

[Describe cross-validation scheme and local validation score.]

\begin{itemize}
    \item \textbf{CV Scheme}: [e.g., Stratified 5-Fold, TimeSeriesSplit]
    \item \textbf{Local CV Score}: [your cross-validation metric]
    \item \textbf{Public LB Score}: [your public leaderboard score]
    \item \textbf{Private LB Score}: [your private leaderboard score]
\end{itemize}

% ==================== POST-MORTEM ANALYSIS ====================
\section{Post-Mortem Analysis}

 [Placeholder: Reflect on what worked, what didn't, and lessons learned.]

\subsection{What Worked}

\begin{itemize}
    \item Successful approach 1 and why it worked
    \item Successful approach 2 and why it worked
    \item Successful approach 3 and why it worked
\end{itemize}

\subsection{What Didn't Work}

\begin{itemize}
    \item Failed approach 1 and why it failed
    \item Failed approach 2 and why it failed
    \item Failed approach 3 and why it failed
\end{itemize}

\subsection{Lessons Learned}

\begin{itemize}
    \item Key takeaway 1
    \item Key takeaway 2
    \item Key takeaway 3
\end{itemize}

\subsection{Future Improvements}

[Placeholder: What would you do differently with more time/resources?]

\begin{itemize}
    \item Improvement 1: e.g., try different model architecture
    \item Improvement 2: e.g., more extensive feature engineering
    \item Improvement 3: e.g., ensemble more diverse models
\end{itemize}

\end{document}
